\documentclass[11pt]{article}
\usepackage{icmlko}
\begin{document}

\title{판의 여덟 중 하나는 갈래를 맞히는 것이다 --- 비슷한 것들 사이에서
챔피언은 $0.44$ 다}
\author{월드모델 실험실}
\date{노트 470 · 2026년 8월 3일}
\maketitle

\begin{abstract}
노트 468 이 집 밖에서 \emph{``갈래 안''과 ``수집 전체''가 부호를 바꾼다}는
것을 보였고, 노트 469 가 \emph{판의 단위는 도메인 전체}임을 코드로
확인했다. 그러면 판 자신에게 물어야 한다 --- \textbf{판 점수의 얼마가
무리 \emph{사이} 순위인가.} 배포에서 IP 를 고를 때는 대개 \emph{비슷한
것들 사이}에서 고르기 때문이다.

도메인마다 유보 행을 갈래(\texttt{gen})로, 없으면 표지(\texttt{grp})로
갈라 무리별 rho 를 크기 가중평균했다.

\begin{center}
\begin{tabular}{lrrr}
\hline
 & 여섯 도메인 가중평균 & \\
\hline
판 (전체 --- 지금 쓰는 값) & $\mathbf{0.4969}$ & \\
무작위 같은 크기로 쪼갬 & $0.4932$ & 크기 몫 $-0.0037$ \\
\textbf{무리 안} & $\mathbf{0.4377}$ & \textbf{동질성 몫 $-0.0555$} \\
\hline
\end{tabular}
\end{center}

\textbf{크기는 아무 상관이 없다.} 무리와 \emph{똑같은 크기 분포}로
무작위로 쪼개면 $0.4932$ 로 전체와 거의 같다. 떨어지는 것은 오직
\textbf{무리가 동질적일 때}다.

\textbf{판 $0.4969$ 중 $0.056$, 즉 여덟 중 하나가 무리 사이 순위다.}
비슷한 것들 사이에서 고르면 챔피언은 $\mathbf{0.4377}$ 이다.
\end{abstract}

\begin{figure}[t]
\centering
\includegraphics[width=\textwidth]{figs/f1.pdf}
\caption{\textbf{가} --- 도메인마다 셋. 파랑(무작위 같은 크기)이 회색
(전체)에 붙어 있고 빨강(무리 안)만 떨어진다. \textbf{나} --- 손실의
분해.}
\end{figure}

\section{도메인마다 다르고, 차이가 크다}

\begin{center}
\begin{tabular}{lrrr}
\hline
도메인 & 전체 & 무리 안 & 동질성 몫 \\
\hline
펀딩 (\texttt{grp} 10무리) & $+0.3640$ & $+0.2317$ & $\mathbf{-0.1337}$ \\
애니 (\texttt{gen} 6무리) & $+0.5731$ & $+0.4856$ & $-0.0824$ \\
모바일 (\texttt{gen} 8무리) & $+0.6223$ & $+0.5549$ & $-0.0532$ \\
웹툰 (\texttt{grp} 3무리) & $+0.4798$ & $+0.4595$ & $-0.0193$ \\
세계애니 (\texttt{gen} 4무리) & $+0.5616$ & $+0.5452$ & $-0.0146$ \\
\textbf{만화} (\texttt{gen} 4무리) & $+0.3478$ & $+0.3621$ & $\mathbf{+0.0170}$ \\
\hline
\end{tabular}
\end{center}

\textbf{펀딩은 점수의 $37\%$ 가 무리 사이다.} 펀딩의 \texttt{grp} 은
스물한 종이라 사실상 ``어떤 종류의 펀딩인가''를 맞히는 것이고, 그 종류를
빼면 남는 게 $+0.23$ 이다. 반대로 \textbf{만화는 무리 안에서 오히려
오른다} --- 만화의 축은 갈래 \emph{안}에서 더 잘 듣는다.

\emph{판에서 제일 높은 도메인이 제일 많이 부풀려진 것도 아니다} ---
모바일($+0.62$)의 손실은 중간이고 펀딩($+0.36$)이 제일 크다.

\section{이것이 바꾸는 것}

\textbf{판을 안 바꾼다.} 순위는 그대로 유보 표본의 판 rho 다 --- 지표를
바꾸면 지난 2년과 못 견준다(노트 446 에서 같은 이유로 안 바꿨다).
바꾸는 것은 \textbf{같이 적는 것}이다.

\begin{quote}
\textbf{조항} --- 챔피언 점수를 적을 때 \textbf{무리 안 판}을 같이
적는다. 잴 수 있는 도메인만(지금 여섯)이고, 무리는 \texttt{gen} 우선 ·
없으면 \texttt{grp} · 20행 이상. 오늘 값은 \textbf{전체 $0.4969$ · 무리
안 $0.4377$}(여섯 도메인 기준)이다.
\end{quote}

이유는 노트 446 의 날짜 통제 판과 같다. 그때 최근순 축은 날 판을
$+0.0186$ 벌면서 실력을 $-0.0294$ 깎았고, 통제 판을 같이 안 찍었으면
못 봤다. \textbf{무리 사이 순위만 잘 맞히는 변경도 같은 방식으로 판을
속일 수 있다} --- 이제 그것도 찍힌다.

\section*{남은 것}

\begin{itemize}
\item \textbf{어느 쪽이 배포인가} --- 여전히 사람에게 물을 것이다. 다만
이제 두 수를 다 들고 물을 수 있다: \emph{``전체에서 고르면 $0.50$,
비슷한 것들 사이에서 고르면 $0.44$ 입니다.''}
\item \textbf{펀딩의 grp 스물한 종} --- 그 표지가 무엇인지 다시 봐야
한다. 점수의 $37\%$ 를 그것 하나가 낸다.
\item \textbf{다섯 도메인은 못 쟀다}(게임 · 도서 · 시장팝업 · 아이돌 ·
팝업) --- 무리 열쇠가 없거나 무리가 둘 미만이다. 시장팝업은 노트 454 가
보인 대로 애초에 칸이 열여덟뿐이다.
\end{itemize}

\end{document}
